\section{Spiritual Warfare, Exorcism, and the Masonic Name}

The Church's condemnation of Masonic membership is not exhausted by sociology.  A baptized person really has been delivered from the dominion of darkness and transferred into the kingdom of the beloved Son (Col 1:13); he has renounced Satan, professed the triune God, and been incorporated into Christ.  An initiation that places him under a religiously multivalent architect, a second passage from darkness to light, and progressively disclosed obligations is therefore not merely an imprudent club affiliation.  It objectively contradicts the baptismal form of Christian life.  In that exact sense the conflict is spiritual.

That conclusion does not authorize a second error: treating every Masonic member, family history, illness, temptation, or misfortune as evidence of extraordinary demonic action.  The devil is a creature, not a rival first principle.  Divine Providence remains sovereign; human beings remain responsible agents; natural, medical, psychological, social, and moral causes retain their integrity.  The Church's rites resist both rationalist denial and occult inflation.  They name the enemy while denying him the power to explain everything.

\begin{massbox}{The reoriented thesis}
The deepest contrast is between two initiations and two renunciations.  Baptism renounces Satan, confesses Father, Son, and Holy Spirit, forgives sin, and incorporates the person into Christ's visible Body.  Masonic initiation constructs a transconfessional ritual fraternity whose common form cannot be governed by that confession.  Catholic spiritual warfare therefore begins in Baptism and continues through faith, prayer, the sacraments, virtue, and obedience to the Church.  Exorcism is a subordinate sacramental ministry exercised under ecclesiastical authority---not a secret counter-initiation, a historical shortcut, or a universal diagnosis of Masons and their descendants.
\end{massbox}

\subsection{Five genres hidden by the phrase ``an exorcism prayer''}

Catholic sources use several kinds of prayer against evil.  They differ in minister, authority, object, and evidentiary force.  Collapsing them is the principal cause of confusion in claims about a ``Masonic exorcism.''

\begin{fourstable}{Genre}{Minister and object}{Ecclesial status}{What it can establish}
Baptismal and catechumenal exorcisms & The Church prays over catechumens and those to be baptized as part of Christian initiation. & Promulgated liturgy; distinct from the major exorcism of a possessed person. & The ordinary Christian grammar: renunciation of Satan, faith in the Trinity, liberation and incorporation in Christ.\\
Major exorcism & A bishop, or a suitable priest with his special and express permission, addresses a person judged possessed. & Reserved sacramental governed by CIC can.~1172 and the Roman Ritual. & The Church's authoritative prayer and discipline in a concrete case; never a general historical proof about an association.\\
Supplication in particular circumstances & Authorized liturgical prayers concerning attacks upon the Church or extraordinary affliction of places or things. & The current ritual's Appendix I; use is governed by the book and competent authority. & The Church can pray corporately against diabolic action without naming Freemasonry or diagnosing every opponent.\\
Private deliverance or renunciation prayer & A person petitions God, rejects sinful ties, or, where authority permits, a cleric prays for liberation. & May be sound private prayer; some books have publication approval.  It is not thereby the Roman Ritual. & A pastoral author's proposal and the devotion's reception, not universal doctrine or independent proof of embedded claims.\\
Unofficial internet formula & Minister, text, edition, and authority often vary; labels such as ``exorcist only'' may be supplied by a ministry site. & No liturgical status follows from online circulation. & At most, evidence that a formula circulates in a contemporary ministry network.\\
\end{fourstable}

The distinction between deprecative and imperative language cuts across these genres.  A deprecative prayer asks God to deliver; an imperative formula directly commands an evil spirit.  The latter is not a spiritual technology made safe by copying words.  Authority, vocation, the person prayed over, and the Church's law matter.  The CDF's 1985 letter recalls that the faithful may not use the formula against Satan and the fallen angels published by order of Leo XIII, still less the complete exorcism.  The current Roman discipline reserves the major rite and forbids improvising it.  A webpage heading cannot confer the jurisdiction that canon law requires.

\subsection{Leo XIII's 1890 exorcism: official, but not Masonic by name}

The most important historical control is the official record.  The \emph{Acta Sanctae Sedis} for 1890--91 prints an \emph{Exorcismus in Satanam et angelos apostaticos}, ``published by order of Leo XIII.''  The formula invokes Christ's authority, the intercession of Mary and Michael, and the Church's divine commission.  In its imperative portion it addresses

\begin{quote}
\latin{omnis congregatio et secta diabolica}
\end{quote}

---``every diabolical congregation and sect.''  It prays against snares laid for the Church and against assaults upon her liberty.  The accompanying rescript, dated 18 May 1890, attached an indulgence for bishops and priests who recited it while legitimately authorized by their ordinaries.

The text does \emph{not} contain \emph{masson}, \emph{massonicus}, ``Masonry,'' ``Freemasonry,'' or the name of a Masonic obedience.  Its broad language can readily be read in the same spiritual and political world as Leo's anti-Masonic teaching: \emph{Humanum genus} preceded it in 1884, and \emph{Dall'alto} followed in October 1890.  Chronological and conceptual adjacency, however, is not textual identity.  The official witness does not warrant calling it ``Leo XIII's exorcism of Freemasonry,'' claiming that it was promulgated because of a revealed vision of a lodge, or using \emph{secta diabolica} as though it were an encoded proper name.

It must also be distinguished from the short prayer to St.~Michael ordered among the prayers after Low Mass.  The familiar short prayer and the longer exorcism have related Leonine reception histories but are different texts with different liturgical and disciplinary histories.  The current \emph{De Exorcismis et Supplicationibus Quibusdam}, publicly presented in 1999, replaces the former chapter of the Roman Ritual.  Its Appendix I supplies a supplication and exorcism for particular circumstances of the Church analogous to the Leonine material; the USCCB's official presentation of the rite does not identify it as a Masonry-specific formula.

\subsection{The modern prayer family that actually names Freemasonry}

The best-known degree-by-degree prayer is modern and has a recoverable Protestant antecedent.  Selwyn R. Stevens published \emph{Prayer of Release for Freemasons \& Their Descendants}; the text itself carries a first copyright date of 1994 and later revision dates.  Public witnesses from 2004 onward and the National Library of New Zealand's record for the expanded seventh edition identify Stevens and Jubilee Resources International.  The prayer walks through alleged obligations and penalties of successive degrees, renounces named spiritual powers, rejects effects attributed to ancestors, and associates particular degrees with bodily or psychological afflictions.

Fr.~Chad Ripperger later adapted this textual family for a Catholic deliverance-prayer collection.  In a 2021 recorded interview about \emph{Deliverance Prayers: For Use by the Laity}, he stated that the Freemasonry prayer had been drafted by a Protestant and that he had ``Catholicized it,'' removed Protestant expressions, changed it, and added material.  Bibliographic records place the first edition of that lay collection in December 2016.  The publisher states that the book carries an imprimatur from the Archdiocese of Denver; accessible front-matter records identify an edition approved by Archbishop Samuel J. Aquila on 23 May 2018.

That history corrects three widespread assertions.  The long prayer is not an ancient Roman exorcism, is not a lost composition of Leo XIII, and did not arise independently as an immemorial Catholic formula.  It is a modern Catholic adaptation of an earlier Protestant deliverance text.  The adaptation may contain distinctively Catholic petitions and may be used pastorally in the circumstances intended by its compiler; its provenance must nevertheless remain visible.

Current ministry websites distribute related texts under several titles:

\begin{itemize}
\item \emph{Prayer to Break the Freemasonic Curse}, a long, degree-by-degree composite assigned by the St.~Michael Center for Spiritual Renewal to Catholic exorcists;
\item \emph{Prayer to Break the Freemasonic Curse---Short Form}, an abridgment that directs readers toward the Archdiocese of Manila Office of Exorcism;
\item \emph{Ratification Prayer for Freemasonry and Other Curses}, assigned to priests and attributed on the page only to ``a Dominican Friar''; and
\item lay or assisted-renunciation versions derived from the Stevens--Ripperger family, sometimes accompanied by instructions to recite the prayer before the Blessed Sacrament or with a priest.
\end{itemize}

The St.~Michael Center is a Catholic nonprofit ministry led by Msgr.~Stephen Rossetti; the Manila office is an archdiocesan ministry.  Those facts deserve more weight than an anonymous repost.  Neither fact makes every hosted text a universally promulgated liturgical book.  The long, short, and ratification pages do not provide an ancient manuscript chain, a decree inserting the formula into the Roman Ritual, or an edition-level approval notice for every wording shown.  The anonymous attribution of the ratification prayer is not enough to establish authorship, date, or approval.

These prayers are included through their identifiable titles, authors or attributions, earliest witnesses, intended ministers, authority status, content architecture, and stable source links.  The complete modern formulas are not reprinted.  They are copyrighted and repeatedly revised; some contain imperative commands assigned to restricted ministers; and detaching a web transcription from its exact edition would defeat the very provenance analysis this section requires.  A historical-theological article should enable the reader to identify and evaluate the texts without becoming an unauthorized exorcism manual.

\subsection{Provenance and authority matrix}

\begin{longtable}{@{}>{\RaggedRight\arraybackslash}p{.16\textwidth}>{\RaggedRight\arraybackslash}p{.15\textwidth}>{\RaggedRight\arraybackslash}p{.18\textwidth}>{\RaggedRight\arraybackslash}p{.18\textwidth}>{\RaggedRight\arraybackslash}p{.15\textwidth}@{}}
\toprule
\textbf{Text or family} & \textbf{Earliest verified witness} & \textbf{Names Masonry?} & \textbf{Status and intended user} & \textbf{Evidentiary use}\\
\midrule
\endfirsthead
\toprule
\textbf{Text or family} & \textbf{Earliest verified witness} & \textbf{Names Masonry?} & \textbf{Status and intended user} & \textbf{Evidentiary use}\\
\midrule
\endhead
\emph{Exorcismus in Satanam et angelos apostaticos} & \emph{Acta Sanctae Sedis} 23 (1890--91), rescript of 18 May 1890 & No; it says ``every diabolical congregation and sect.'' & Official Leonine formula; the rescript concerned bishops and duly authorized priests; later discipline restricts use. & Official Catholic demonology and prayer for the Church; not proof of a named Masonic rite.\\
Current Roman Ritual and Appendices I--II & Revised rite publicly presented 26 Jan. 1999; vernacular editions approved later & No Masonry-specific formula located in the official public inventory consulted. & Major rite reserved; Appendix I governed liturgically; Appendix II supplies private supplications. & Controlling present ritual categories and discipline.\\
Stevens, \emph{Prayer of Release for Freemasons \& Their Descendants} & Textual copyright begins 1994; revised public copies and later library record & Yes, explicitly and degree by degree. & Protestant or evangelical deliverance and renunciation composition; repeatedly revised. & Earliest verified ancestor of the modern long prayer family; not Catholic liturgy.\\
Ripperger, \emph{Prayers to Break the Freemasonic Curse} & Lay collection first published Dec. 2016; 2018 approved edition; compiler's 2021 account of adaptation & Yes. & Catholicized private deliverance formula in a prayer book advertised with a Denver imprimatur. & Evidence of a modern Catholic pastoral adaptation; contingent assertions require separate proof.\\
St.~Michael Center long and short forms & Contemporary online ministry publication; short page links to Manila office & Yes. & Ministry texts assigned to specified users; not shown to be a Roman liturgical edition. & Evidence of current reception and abbreviation, not an independent ancient witness.\\
\emph{Ratification Prayer for Freemasonry and Other Curses} & Contemporary St.~Michael Center webpage & Yes. & Priest-facing prayer attributed only to an unnamed Dominican friar. & Present pastoral usage only; authorship, date, textual genealogy, and formal approval remain unverified.\\
\bottomrule
\end{longtable}

\subsection{What an imprimatur says---and does not say}

Canon 826 \S3 requires the local ordinary's permission for publication of prayer books.  Under canons 829--830, ecclesiastical approval concerns the submitted original text or edition and the censor's judgment that it contains nothing contrary to faith or morals.  It does not automatically extend to a new translation, altered edition, excerpt, or website transcription.  Nor does it make a private prayer part of the Roman Ritual, define its historical hypotheses, canonize its author's pastoral diagnoses, or authorize a reader to perform a reserved exorcism.

This matters especially because the online variants are not textually inert.  Long and short forms omit, add, or rearrange clauses; headings assign them to different ministers; derivative pages may not identify the exact approved edition.  ``Has an imprimatur,'' ``is prayed by an exorcist,'' ``appears on a Catholic ministry site,'' and ``is an official exorcism of the Church'' are four different claims.  Only the first three are established for particular witnesses in this modern family, and not all three for every witness.

\subsection{Prayer language is not documentary proof}

The long modern formula combines several kinds of proposition: Christian renunciation of sin; petitions to Christ; rejection of Masonic oaths; names or interpretations assigned to degrees; claims about demons; claims about curses passing through ancestry; and mappings between degrees and medical, anatomical, psychological, or social effects.  A prayer can faithfully ask deliverance while containing prudential or factual applications that do not authenticate themselves by being prayed.

Thus a clause naming Baphomet, Asmodeus, witchcraft, Lucifer, a particular bodily organ, or a numbered degree is not evidence that every lodge invokes that spirit, that the degree historically bears that demon's name, or that a descendant's disease was caused by the ancestor's enrollment.  Such propositions require ritual texts, institutional records, medical evidence, or disciplined case discernment appropriate to the claim.  Statements attributed to demons in an exorcism cannot function as a historical archive: the devil is a liar, the Church's rites do not exist to satisfy curiosity, and sensational interrogation is explicitly rejected by Catholic discipline.

This rule strengthens rather than secularizes the article.  Spiritual reality is affirmed on the authority of Revelation and the Church, not on the prestige of an unverifiable anecdote.  Historical claims are established historically.  Medical conditions are evaluated medically.  Extraordinary affliction is discerned under ecclesiastical authority.  Human culpability is judged according to object, knowledge, and consent.  No one level is allowed to impersonate all the others.

\subsection{Ancestors, consequences, and the grace of Baptism}

The phrase ``Masonic generational curse'' needs especially careful treatment.  Scripture recognizes that sin harms families and societies across time: children can inherit poverty, trauma, learned vice, scandal, broken trust, bodily risk, and unjust structures.  Aquinas likewise distinguishes personal guilt from temporal effects that reach others through natural and social bonds (ST I--II, q.~87, a.~8).  None of this makes a descendant personally guilty of an ancestor's sin.

The universal sacramental doctrine is clear.  Original sin is transmitted in a unique manner but is not personal fault in Adam's descendants (CCC 405).  Baptism forgives all sin and all punishment due to sin; temporal weakness and consequences remain for struggle and growth (CCC 1263--1264).  Ezekiel rejects imputing a father's guilt to his son (Ezek 18:20), and Christ refuses a simple inherited-guilt diagnosis of the man born blind (John 9:1--3).

The Spanish Episcopal Conference's 2024 doctrinal note on intergenerational healing applies these principles to a modern pastoral theory that diagnoses present illness from ancestors' unforgiven sins and proposes liberation by special prayer, exorcism, or Mass.  The note traces that theory to twentieth-century authors, judges it foreign to Catholic tradition, rejects the transmission of personal guilt, warns that it compromises Baptism and Reconciliation, and still acknowledges real familial, social, psychological, and bodily consequences.  This is a particular episcopal intervention, not a new universal dogmatic definition; its argument rests on universal sacramental teaching and directly tests the strongest claims in the modern Freemasonry-prayer family.

Catholic theology need not declare it metaphysically impossible that another person's curse or occult act could occasion a specific external affliction.  It does reject automatic lineage diagnosis, personal inherited guilt, magical causal chains, and the presumption that a disease pattern reveals a Masonic ancestor.  A case-specific claim needs case-specific discernment.  The descendant's ordinary path is not fear-driven excavation of a hidden family degree but life in Christ: Baptism, repentance for one's own sins, forgiveness of others, sacramental confession where needed, Eucharistic communion, prayer, virtue, and prudent care.

\subsection{The pastoral order of operations}

For an active Mason, the first questions are public and moral: Is he enrolled?  What rites and obligations has he accepted?  Will he cease participation, resign, renounce incompatible promises, confess grave sin, and repair scandal or injustice?  Those actions address the objective conflict.  They do not require a prior diagnosis of possession or an inventory of alleged degree-demons.

For a relative or descendant who never joined, no Masonic guilt is his to confess.  He may reject any family glorification of the lodge, forgive injuries, disclose concrete occult practices or grave sins of his own, and ask God for protection.  He should not be taught that he is secretly enrolled by blood.  Unexplained medical or psychological suffering warrants appropriate professional care; suspected extraordinary phenomena should be taken to the pastor and diocesan protocol rather than tested with internet formulas.

For clergy, the order is equally ecclesial: ordinary pastoral care, sacramental reconciliation, sound catechesis, factual withdrawal, safeguarding, medical and psychological evaluation where indicated, and referral to the bishop's designated minister.  A priest without the required permission does not acquire it from the gravity of the story.  A layperson may pray to Christ for protection and deliverance, renew baptismal promises, use approved private supplications proper to the laity, and invoke Mary, Michael, and the saints; he must not appropriate the major rite or the restricted Leonine formula.

\begin{studybox}{What the prayer evidence finally proves}
The official Roman tradition proves that the Church takes demonic opposition seriously and governs exorcism under sacramental and episcopal authority.  Leo XIII's 1890 text proves that he prayed against diabolical congregations attacking the Church; it does not name Masonry.  The Masonry-specific texts prove that a modern Protestant deliverance formula was adapted and received in some Catholic exorcist and deliverance ministries.  They do not, by themselves, prove their degree-by-degree history, demonology, medical mappings, or automatic ancestral transmission.  The strongest Catholic response is consequently more, not less, supernatural: the victory of Christ given in Baptism, ordinarily lived through the Church's sacraments and virtues, with extraordinary ministry kept obedient to the authority Christ gave his Church.
\end{studybox}
